% Automatisch aus der Word-/LaTeX-Konvertierung {\"u}bernommen.
\chapter{Validit{\"a}t}


\section{Interne Validit{\"a}t}

\section{Externe Validit{\"a}t}
Bei der externen Validit{\"a}t von einer Untersuchung handelt es sich eine Beschreibung davon, wie weit sich die Ergebnisse auf andere Bedingungen als die untersuchten, {\"u}bertragen lassen. Dabei ergeben sich f{\"u}r diese Arbeit mehrere Einschr{\"a}nkungen der {\"U}bertragbarkeit.

Erstens wurden alle Messungen auf genau einer Hardware-Plattform durchgef{\"u}hrt. Dabei handelt es sich um den nRF52840 DK mit einem 32-Bit-Arm-Cortex-M4F-Prozessor bei 64 MHz Taktfrequenz (siehe 4.2). Andere Mikrocontroller-Architekturen, wie bspw. Cortex-M0+, dem Cortex-M7 oder RISC-V-basierte Systeme haben andere Pipeline-Tiefen, Speicherzugriffszeiten und Interrupt-Controller-Architekturen. Die absoluten Zyklenzahlen dieser Arbeit k{\"o}nnen daher nicht direkt auf andere Hardware {\"u}bertragen werden. Ob sich dabei die relativen Unterschiede zwischen den Konfigurationsstufen auch auf andere Hardware {\"u}bertragen lassen, wurde nicht untersucht.

Zweitens wurde nur Zephyr in der Version v3.7.2 (LTS) mit dem Zephyr SDK 0.16.9 (GCC 12.2.0) untersucht (siehe 4.6.2, 4.6.4). Der Kernel und auch einzelne APIs k{\"o}nnen sich zwischen den Zephyr-Versionen {\"a}ndern.

Drittens wurden nur drei Konfigurationsstufen aus dem gr{\"o}{\ss}eren Kconfig-Raum von Zephyr untersucht (siehe 4.4). Andere Funktionskombinationen, wie z.B. Bluetooth, USB, Dateisysteme, Userspace/ MMU-Isolation, wurden nicht gemessen und ihre Kosten lassen sich auch nicht aus den ergebenen Daten ableiten.

Viertens wurde als praxisnahes Anwendungsbeispiel nur ein IoT-Sensor-Szenario {\"u}ber Thread, CoAP und DTLS untersucht (siehe 6.4). Andere verbreitete IoT-Kommunikationsstacks, wie bspw. Bluetooth Low Energy, Wi-Fi oder Matter {\"u}ber Wi-Fi/ Ethernet, wurden nicht untersucht und k{\"o}nnen daher andere Ressourcenkosten aufweisen.

In Kapitel 8.4 wurde eine Einordnung dieser Arbeit gegen Silva et al. (2019) durchgef{\"u}hrt, welche einen externen Referenzpunkt bringt. Wenn die eigenen Fu{\ss}abdruck-Werte mit denen einer unabh{\"a}ngigen Studie {\"u}bereinstimmen oder sich {\"a}hneln, dann wird das Vertrauen in die {\"U}bertragbarkeit der Methodik dadurch verst{\"a}rkt. Es ersetzt jedoch auf der anderen Seite keine eigene Messung auf anderer Hardware oder mit anderen Zephyr-Versionen.

\section{Messbedingter Zusatzaufwand}
Jede Zeitmessung an einem laufenden System beeinflusst das gemessene Ergebnis durch ihre eigene Durchf{\"u}hrung. Dieser Effekt wird Messaufwand genannt. Der Messaufwand in der Messmethodik, die f{\"u}r diese Arbeit entwickelt wurde, ist vor allem das Auslesen des DWT-Zyklenz{\"a}hlers (siehe 6.2.3). Es beansprucht n{\"a}mlich vor und nach dem zu messenden Vorgang einen kleinen Zeitanteil.

Um diesen Anteil aus den Ergebnissen rauszurechnen, wird der Messaufwand vor jeder Messreihe kalibriert. Der Zyklenz{\"a}hler wird 100-mal direkt hintereinander ausgelesen, ohne dass dazwischen irgendwas gemessen wird. Dann wird das Minimum aller 100 Differenzen als Kalibrierungswert genommen. Dabei ist bewusst das Minimum ausgesucht worden. St{\"o}reinfl{\"u}sse k{\"o}nnen n{\"a}mlich eine einzelne Kalibrierungsmessung nur verl{\"a}ngern, nicht verk{\"u}rzen. Dadurch wird der echte Grundaufwand durch das Minimum so gut wie m{\"o}glich angen{\"a}hert.

Dieser kalibrierte Wert wird dann anschlie{\ss}end von jedem einzelnen rohen Messwert abgezogen, bevor er in die laufende Statistik (siehe 6.2.4) einflie{\ss}t. In den Messungen lag der kalibrierte Messaufwand bei 12 Zyklen. In den dokumentierten Kalibrierungen wurde derselbe Kalibrierungswert beobachtet. Daraus folgt keine allgemeine Unabh{\"a}ngigkeit des Messaufwands von der Konfiguration.

Die Kalibrierung erfasst jedoch nur den Zeitaufwand des Z{\"a}hler-Auslesens selbst. Daher werden nicht andere m{\"o}gliche Seiteneffekte, die durch den Code entstehen k{\"o}nnen, wie bspw. ver{\"a}nderte Flash-Adressen von benachbarten Funktionen erfasst (siehe 9.1). Wie stark solche {\"A}nderungen des Flash-Layouts die einzelnen Messwerte beeinflussen, wurde in dieser Untersuchung nicht separat gemessen. Sie bleiben daher eine m{\"o}gliche Quelle von systematischen Abweichungen zwischen den Firmware-Builds.

Der Selbsttest (siehe Kapitel 5) ist eine weitere Kontrolle f{\"u}r die Nachvollziehbarkeit der gemeinsam benutzten Zeitmessungs- und Auswertungskette. Er pr{\"u}ft die gemeinsame Wirkung von Zeitbasis, Abzug des Messaufwands und Auswertung anhand einer vorgegebenen Busy-Wait-Dauer auf Nachvollziehbarkeit. Die einzelnen Teile werden dadurch nicht unabh{\"a}ngig voneinander validiert. Das hierbei definierte Abnahmekriterium ist in allen Selbsttests durchgef{\"u}hrt und erf{\"u}llt worden.

\section{Limitierungen}
% ===== END KONVERTIERTER INHALT =====
